% Section 1 — Introduction (charter §4.1; issue T4-1 prose, drawing on T1-2/T1-3)
\section{Introduction}
\label{sec:intro}

Does the \emph{format} in which a program is presented to a code language model
change what the model can predict about it? The question is usually asked
empirically --- pick a family of representations, hold the model fixed, and
measure the score --- and the empirical companion to this paper,
\texttt{preprint-v1}, asked it exactly that way on the Lumen apparatus: a frozen
model queried under a ladder of code representations C2--C4 against an
information-matched control C1${}^{+}$. Its central, pre-registered cell
returned a clean null (a rank-biserial effect $r_{\mathrm{rb}}=+0.10$ at raw
$p=0.36$, no rubric ceiling). This paper is the theoretical sibling of that
record. It does not revise the preprint, re-run its scorers, or touch its frozen
analysis; it asks \emph{why} the null is the outcome one should have expected,
and --- the harder half --- where an effect, if one exists at all, is permitted
to live.

The answer we defend has two halves, and the paper is honest about which half
each result belongs to. The first is negative and is a theorem. Once
``representation'' and ``predictor'' are made into measurable objects on a common
probability space (\cref{sec:prelim}), a representation is a measurable map that
can only \emph{coarsen} the input $\sigma$-algebra, never enrich it. The
data-processing inequality (\cref{lem:dpi}) then forbids any representation from
creating predictive information about the target, and at the Bayes optimum this
sharpens to an equality: under \emph{information parity} --- two representations
retaining the same mutual information with $Y$ --- the Bayes-optimal risk is
\emph{format-invariant} (\cref{thm:irrelevance}). At constant information the
optimum cannot tell the formats apart. The Lumen T1 null is, on this reading,
not a weak or disappointing positive but the \emph{predicted} outcome of the
information-limit regime, and treating it as anything else is the overclaim our
honesty guardrails exist to prevent.

The second half is positive, and it is deliberately \emph{not} a claim that any
effect exists. A frozen LLM is not the Bayes posterior; it is a
finite-capacity predictor that does not attain the information limit, and the
gap between what it achieves and that limit is precisely where format is free to
act. We bound that gap. For bounded learners the achievable
parity-preserving advantage is at most a capacity budget $\Phi$ built from an
approximation (expressivity) gap and a Rademacher estimation term, and carrying
\emph{no} mutual-information term whatsoever (\cref{thm:bounded-gap}); the bound
is one-sided and sign-agnostic --- it says how large an effect \emph{can} be,
never that one exists or which format wins (\cref{cor:size-gap}). Because the
budget is non-empty only in a named finite-resource, capacity-separated window
(\cref{prop:regime}), the theory does something a bare null cannot: it
\emph{localizes} the regime in which a representation effect is permitted,
forbids it everywhere else, and turns that localization into a list of
falsifiable predictions on the same apparatus (\cref{sec:predictions}).

\subsection{The format-versus-information question, stated formally}
\label{subsec:intro-question}

Informally the empirical question conflates two things that the theory must
separate: whether a representation carries different \emph{information} about the
target, and whether, at fixed information, a particular \emph{learner} extracts
different value from it. The Lumen C1${}^{+}$ control exists precisely to fix the
first so the second can be asked alone, and the formal apparatus of this paper
exists to say what ``fix the first'' means exactly. We write $\rep:\Xspace\to
\Zspace_{\rep}$ for a representation and ask of a pair $\rep,\rep'$:
\begin{enumerate}
  \item[(Q1)] \emph{Information.} Do $\rep$ and $\rep'$ retain the same
    predictive content, $\MI{Y}{\rep(X)}=\MI{Y}{\rep'(X)}$? This is
    \emph{information parity} (\cref{def:parity}); its categorical form is a
    retraction (\cref{sec:category}) and its quantitative failure is a transport
    defect (\cref{sec:parity-defect}).
  \item[(Q2)] \emph{Value to a bounded learner.} Given parity, can a
    finite-capacity predictor still do measurably better under one than the
    other? This is the representation advantage $\advantage(\rep,\rep')$
    (\cref{def:advantage}), and the bounded half of the paper is the statement
    that whatever value exists here is capacity, not information
    (\cref{sec:bounded}).
\end{enumerate}
Phrased this way the empirical null is the answer to a \emph{composite} of Q1 and
Q2: the C4-vs-C1${}^{+}$ comparison was engineered so Q1 is answered ``same''
(by construction), leaving the recorded effect to measure Q2 alone. The theorem
of \cref{sec:parity} is that the Q1 channel is genuinely closed at the optimum;
the bound of \cref{sec:bounded} is that the Q2 channel, while open, is narrow and
capacity-controlled. The negative half disposes of the information explanation;
the positive half says where to hunt for the capacity one.

\subsection{Contributions}
\label{subsec:intro-contributions}

\begin{enumerate}
  \item \textbf{A measure-theoretic frame in which the question is exact}
    (\cref{sec:prelim}). Representations become measurable maps and predictors
    become Markov kernels on a shared standard Borel space, so that ``the
    predictor only sees the representation'' is the statement that its kernel is
    $\repsig$-measurable, and the data-processing and sufficiency arguments
    become exact statements about conditional expectations rather than heuristics
    about ``information.''
  \item \textbf{An irrelevance theorem and its rate--distortion corollary}
    (\cref{sec:parity}). Under information parity the Bayes risk is
    format-invariant (\cref{thm:irrelevance}); when parity fails, the inflation
    of Bayes risk is controlled by the information lost --- exactly the lost bits
    under log loss, and a $\sqrt{\cdot}$ envelope for any bounded loss
    (\cref{cor:ratedist-gap}). This is the formal content of the negative half
    and the disciplined statement of guardrail~G1.
  \item \textbf{A bounded-learner advantage bound} (\cref{sec:bounded}). At
    constant information the achievable advantage is capped by a capacity budget
    $\Phi$ with no mutual-information term (\cref{thm:bounded-gap}), two-sided in
    magnitude (\cref{cor:size-gap}), and non-vacuous only in a named
    finite-resource window (\cref{prop:regime}). This is the positive half,
    stated as a bound and never as an existence claim (guardrail~G2).
  \item \textbf{Two go/no-go refinements that earned their place}
    (\cref{sec:category}, \cref{sec:parity-defect}). A categorical reading shows
    operational parity is a retraction whose round-trip defect is invisible to
    both the mutual information and the Bayes risk yet real
    (\cref{thm:retraction}); a transport/information-geometry reading metrizes
    that defect and proves it dominates the swap gap of any fixed bounded
    predictor (\cref{thm:defect-risk-gap}). Each survives only because it proves
    something the core two sections cannot (charter §2).
  \item \textbf{Retrodiction without overfitting, and falsifiable predictions}
    (\cref{sec:anchoring}, \cref{sec:predictions}). We map each of the three
    frozen Lumen cells to the regime it occupies --- T1 to the information limit,
    T3 to sub-frontier saturation, T2 to a permitted-but-unidentified bounded
    effect --- claiming only consistency, never confirmation; and we discharge
    the standing charge that a math-dressed null is unfalsifiable by issuing five
    sign-aware, falsifier-bearing predictions on the Lumen knobs
    (\cref{tab:pred-summary}).
\end{enumerate}

\subsection{Relationship to the empirical preprint, and what this paper is not}
\label{subsec:intro-companion}

\texttt{preprint-v1} is the frozen empirical record; this is its theoretical
companion, and the two are kept deliberately at arm's length. We import three
numbers from the preprint and produce none: the T1 null
($r_{\mathrm{rb}}=+0.10$, $p=0.36$), the T3 ceiling ($r_{\mathrm{rb}}=+0.43$ with
$93.7\%$ of scores saturated, effective $\nsample\approx5$), and the
Holm-absorbed T2 positive ($r_{\mathrm{rb}}=+0.58$, raw $p=0.023$, non-rejection
over the nine-cell family). Every appearance of these in \cref{sec:anchoring} is
read off the record; none is a new measurement. The theory's own contribution is
the structure around them --- the theorems, and the predictions that put the
theory at risk of being wrong.

Three disclaimers bound the whole paper and are stated once. First, the
irrelevance theorem holds at two \emph{idealizations} neither met by a frozen
LLM --- Bayes-optimality and exact parity --- and every downstream invocation
discloses which it relies on (guardrail~G3, the \texttt{\textbackslash
idealization} boxes throughout). Second, the bounded-learner objects
($\nsample$, $\hypclass$, the ERM $\widehat f_{\rep}$) are read onto a single
fixed model by analogy, so the retrodiction asserts bounds and regimes, never
point values. Third, the two go/no-go sections are provisional by charter and
carry their decision in-source; they are kept only because each earned a theorem.
What this paper is not is a claim that format does not matter, that any effect is
zero, or that the preprint should have found a positive: it is the statement that
\emph{no information-channel effect is expected}, that any real effect is a
bounded-computation phenomenon of capped size, and that the theory names where to
go looking.
